%!TeX root=./tese.tex
% Conclusions. Numbers quoted here are established in the body chapters
% and ledgered in progress/results-ledger.md.
\chapter{Conclusions}
\label{cap:conclusion}

This thesis asked how the pools of a modality-disaggregated
multimodal serving system should be scaled, and defended a specific
answer: each pool on its own demand signal --- the vision pool
forecast-anticipated (its signal forecastable but not
rate-visible), the language pool tail-aware converted. The four
research questions of Chapter~\ref{cap:intro} can now be answered from
evidence.

\paragraph{RQ1 --- Signal sufficiency: no.} The request-rate signal a
conventional autoscaler watches does not carry the vision pool's demand:
the two pools' demands decouple statistically and saturate an order of
magnitude apart on steady traffic, with the encoder pool failing its
tail bound at the trace's own recorded rate --- invisibly to rate- and
utilization-based monitors. Correcting the missizing needs not just more
replicas but a fleet-shape change (worker count and tensor-parallel
shape) that no replica scaler can perform.

\paragraph{RQ2 --- Forecastability: yes, zero-shot.} Three pretrained
time-series foundation models beat every classical baseline in all
twenty (resolution, signal) cells under one rolling-origin protocol;
accuracy degrades only marginally across a twelve-step horizon covering
the median burst onset, covariates are nearly irrelevant at
provisioning cadence, and the p90 band edge is empirically calibrated
(coverage 0.88--0.90) --- what later converts forecasts into
provisionable headroom.

\paragraph{RQ3 --- Platform validity: yes, with measured fidelity.} The
encoder cost model was collected on real silicon for this thesis (and
contributed upstream), and a composed forward matches vLLM's own vision
tower to within $12.8\%$ --- the earlier ${\approx}1.4\text{--}1.6\times$
gap being HuggingFace execution overhead, not model error
(Appendix~\ref{app:engine-sensitivity}). The language pool runs under
calibrated engine timing (a Dynamo discrete-event simulation on AIC
H100 timing, not real H100) whose delta against the simulator's default
polynomial is quantified, including a burst-window SLA violation the
polynomial misses. The control loop itself is anchored to the live
Dynamo stack: on the fastest-onset burst, the winning policy and the
reactive baseline reproduce their simulated verdicts through the real
control plane with mocked engines --- validating the control loop, not
silicon.

\paragraph{RQ4 --- Closed-loop payoff: yes, with the mechanism
identified.} On the encoder pool, the p90-forecast schedule meets the
tail SLA on every window at every cold-start lead, including the
one-minute-onset burst that defeats an identically informed reactive
controller at every lead. It does so because the p90 band holds
anticipatory headroom that matures within the cold-start lead, standing
capacity up before the spike and leaving the service floor intact. The
forecaster is load-bearing but its \emph{accuracy} is second-order: a
cheap forward forecaster (EWMA) at a calibrated quantile, driven through
the same loop, either misses the fast onsets or over-provisions --- so
it is the calibrated band \emph{shape} that carries the result. On the
language pool --- a same-input \emph{conversion} test, not a second
anticipation experiment --- the forecast passes bursts under realistic
replica leads only when combined with a tail-aware capacity conversion
(a measured belief-to-knee ratio); fed through the planner's native
mean-based conversion it fails like the reactive baseline, so the defect
is the conversion, not loop tuning. Over the covered span the policy
holds the SLA in aggregate and inside every episode at
\textbf{0.62$\times$ the deployable static encoder fleet's provisioned
worker-time} (billed at forecast-issue time under a warm-pool readiness
model, 0.59--0.62$\times$ under strict re-warm; the language fleet
passes at 0.96$\times$ its burst-sized static shape), while the reactive
controller's apparent aggregate pass dissolves under per-episode
grading.

\paragraph{The thesis statement, revisited.} Every clause is now
backed: the pools' demands decouple and saturate far apart (RQ1),
the least visible signal is the most persistent and forecastable
zero-shot at provisioning horizons (RQ2), and scaling the vision pool on
those forecasts --- with the language pool on their tail-aware conversion
--- holds the tail objective through the onsets that defeat reactive control
at every tested lead, at roughly two-thirds of the best static
encoder fleet's GPU-time (RQ4).

\paragraph{What bounds these claims.} The evidence is one seven-day
production trace with an offered-work upper bound on decode demand
(faithful in total volume but temporally displaced, licensing no
pointwise provisioning claim); the inter-pool handoff is an analytic
single-stream bound, so the \SI{588}{\milli\second} floor and the cost
ratio assume an embedding-only handoff or a ${\geq}100$\,Gbps fabric;
and the vision-pool control loop is exercised in simulation against the
measured silicon cost (validated to $12.8\%$;
Appendix~\ref{app:engine-sensitivity}), not yet as a live serving loop.
None reverses the direction of the results --- the encoder cost model in
particular is validated rather than conservatively biased --- but each
scopes their generality (Chapters~\ref{cap:closedloop} and
\ref{cap:twopool} detail them).

\section{Future work}
\label{sec:conc-future}

Three extensions remain. \emph{Completing the physical encoder anchor}:
Chapter~\ref{cap:aiconfigurator} reports a dense whole-tower sweep on the
5090 whose measured $f(N)$ matches the AIC composition across the serving
regime. What remains is one real encoder worker serving a burst replay
end to end and, beyond it, the full closed loop on a real GPU serving
cluster --- the complement to the mocked-engine control-loop anchor of
Chapter~\ref{cap:validation}, and the evaluation modality the systems of
Chapter~\ref{cap:related} already bring. (A structural language-pool
predictor decomposing prefill demand into text-borne and image-driven
components was tested and closed: the decomposition buys nothing even at
onsets, because the image component is least forecastable exactly there;
Chapter~\ref{cap:predictors}.) \emph{The seam as a preview channel}: in a
modality-disaggregated system the vision pool's admitted work is an
exact, short-horizon preview of language-pool arrivals
(horizon $\approx$ encoder sojourn); useless for replica lead, but free
information at scheduling timescales (chunked-prefill planning, KV
pre-allocation), a natural planner extension. \emph{Workload
breadth}: image-size-aware encoder costing when a trace with size
information becomes available, the intermediate model family members,
and fine-tuned forecasters as the forecasting-side upper bound.
